\begin{pastexam}[label={px:gk3-112-ds-01}]{高考三級資訊處理 112 年 第一題}
（申論；25 分）

某一公司有下圖所示的8個優先順序分別為高或低的待執行工作，且將依順序自A至H每間隔一天的時間放入對應的高優先執行佇列（Queue）或低優先執行佇列（Queue），例如A（低）表示A工作將於第一天放入低優先執行佇列，而C（高）表示C工作將於第三天放入高優先執行佇列。此外，執行每個工作所需完成的時間均於工作名稱下顯示，例如執行A工作需要2天時間完成，而執行B工作需要1天時間完成。最後，各個工作的執行規則為，當高優先執行佇列內有工作待完成時，須優先執行該佇列內的工作（由第一個開始執行），直到高優先執行佇列內沒有任何待完成工作時，方可執行低優先執行佇列內的工作（由第一個開始執行）。

\begin{center}
自A至H每間隔一天的時間放入對應的高優先佇列或低優先佇列\\[2pt]
\small\renewcommand{\arraystretch}{1.2}\setlength{\tabcolsep}{4pt}
\begin{tabular}{|*{8}{c|}}
\hline
H（低）& G（高）& F（高）& E（低）& D（高）& C（高）& B（低）& A（低）\\ \hline
1 & 2 & 1 & 1 & 2 & 2 & 1 & 2 \\ \hline
\end{tabular}
\end{center}
\begin{itemize}[leftmargin=2.2em]
\item[㈠] 試計算執行此8個工作需要多少天方可完成。（10分）
\item[㈡] 試計算此8個工作自放入佇列至開始執行的平均等待時間。（15分）
\end{itemize}
\end{pastexam}

\begin{pastsoln}
答案：㈠ 第 $12$ 天結束時全部完成，共 $12$ 天；㈡ 平均等待 $21/8 = 2.625$ 天。

\textbf{解讀}。 第 $k$ 個工作（A 是第 $1$ 個）在第 $k$ 天開始時放入佇列；以「天」為單位，每天開始時若沒有工作
正在執行，就依規則從佇列取一個工作，從當天開始執行。 工作一旦開始就做完（不中斷）。 等待時間＝開始執行的
那一天 $-$ 放入佇列的那一天。 （本題的排程裡，低優先工作只在第 $1$–$2$ 天與第 $10$ 天以後執行，那段期間
沒有高優先工作進來，所以就算規定高優先工作可以中斷低優先工作，結果也一樣。）兩個佇列都是 FIFO 佇列
（資結書第~6.2~節）。

\textbf{逐日追蹤}。 表格的工作順序是由右往左 A 到 H。 佇列內容寫成「工作（剩下天數）」：
\begin{center}
\small
\begin{tabular}{cllll}
\toprule
天 & 放入 & 執行中 & 高優先佇列 & 低優先佇列 \\
\midrule
1 & A（低，2） & A（第 1 天開始） & & \\
2 & B（低，1） & A & & B \\
3 & C（高，2） & C（第 3 天開始） & & B \\
4 & D（高，2） & C & D & B \\
5 & E（低，1） & D（第 5 天開始） & & B, E \\
6 & F（高，1） & D & F & B, E \\
7 & G（高，2） & F（第 7 天開始） & G & B, E \\
8 & H（低，1） & G（第 8 天開始） & & B, E, H \\
9 & & G & & B, E, H \\
10 & & B（第 10 天開始） & & E, H \\
11 & & E（第 11 天開始） & & H \\
12 & & H（第 12 天開始） & & \\
\bottomrule
\end{tabular}
\end{center}
說明幾個關鍵點：第 $3$ 天 A 剛做完，低優先佇列有 B、高優先佇列剛放進 C，高優先先做，所以做 C。 第 $5$、$7$、$8$
天都是高優先佇列有工作，低優先的 B、E、H 一直等。 第 $10$ 天高優先佇列空了，才依 FIFO 做 B、E、H。

\textbf{㈠}。 最後一個工作 H 在第 $12$ 天執行，第 $12$ 天結束時全部完成，需要 $12$ 天。 （檢查：總工作量
$2+1+2+2+1+1+2+1 = 12$ 天，而從第 $1$ 天起處理器每天都在做事，沒有閒置，所以正好 $12$ 天。）

\textbf{㈡}。
\begin{center}
\begin{tabular}{c*{8}{c}c}
\toprule
工作 & A & B & C & D & E & F & G & H & 合計 \\
\midrule
放入（第幾天） & 1 & 2 & 3 & 4 & 5 & 6 & 7 & 8 & \\
開始（第幾天） & 1 & 10 & 3 & 5 & 11 & 7 & 8 & 12 & \\
等待（天） & 0 & 8 & 0 & 1 & 6 & 1 & 1 & 4 & 21 \\
\bottomrule
\end{tabular}
\end{center}
平均等待時間 $= 21 / 8 = 2.625$ 天。 高優先工作平均只等 $(0+1+1+1)/4 = 0.75$ 天，低優先工作平均等
$(0+8+6+4)/4 = 4.5$ 天，這就是兩個優先佇列的效果。
\end{pastsoln}
